\begin{statement}{C.2.UN.33}{Proved}{A properly elliptic normalization component has one of eleven positive-coherent-Euler fibre configurations}
Let $M$ be the smooth minimal model of a normalization
component of a fibre of an unrestricted two-torus degeneration,
and suppose $\kappa(M)=1$. Its connected elliptic Iitaka map
$p:M\to B_0$ is isotrivial, and $g(B_0)\le1$.
Its coherent Euler characteristic belongs to $\{0,1,2\}$.
When it is zero, every fibre of $p$ is a positive multiple
of a smooth elliptic curve. When it is positive,
$B_0\simeq\PP^1$, $q(M)=0$, and the complete configuration
of the fibres of $p$ with singular reduced support is one
of the following rows:
\[
\begin{array}{c|l|c}
\chi(\OO_M)&\text{singular reduced fibres of }p&e(M)\\ \hline
1&\mathrm{II}+\mathrm{II}^*&12\\
1&\mathrm{III}+\mathrm{III}^*&12\\
1&\mathrm{IV}+\mathrm{IV}^*&12\\
1&2\mathrm I_0^*&12\\
1&3\mathrm{IV}&12\\
1&2\mathrm{III}+\mathrm I_0^*&12\\
1&\mathrm{II}+\mathrm{IV}+\mathrm I_0^*&12\\ \hline
2&4\mathrm I_0^*&24\\
2&3\mathrm{IV}^*&24\\
2&2\mathrm{III}^*+\mathrm I_0^*&24\\
2&\mathrm{II}^*+\mathrm{IV}^*+\mathrm I_0^*&24 .
\end{array}                                                   \tag{C.EQ.1}
\]
In each row $p_g(M)=\chi(\OO_M)-1$. All additional
smooth multiple fibres are retained as part of the specified
surface; the table lists its nonsmooth reduced fibres.
This is a necessary classification of the displayed
invariant for actual components. The statement does not
assert occurrence of every abstract surface with these data.
The map $p$ is the elliptic map on the component surface;
the original threefold map has the degeneration disc as target.
\end{statement}

\begin{proof}
Take the actual cyclic component cover of C.2.UN.16.
Pluricanonical pullback gives Kodaira dimension at least
one, and UN15 excludes dimension two. Its minimal model
$E$ therefore has Kodaira dimension one. The results
UN22 and UN29 give its connected elliptic Iitaka map
\[
 q:E\longrightarrow B,\qquad g(B)\le1,                 \tag{C.EQ.2}
\]
with a fixed elliptic fibre, identity integral monodromy,
and only smooth multiple fibres. The cyclic stabilizer
$G$ preserves $q$, because it acts on the pluricanonical
field. The normalization of the original component is
birational to the actual finite quotient $E/G$, through
the equivariant graph and minimal-model comparisons of
UN19--20. Thus its relatively minimal elliptic model is
the model of that quotient over $B/G$.

Let $H$ be the kernel of $G\to\operatorname{Aut}(B)$.
On a general elliptic fibre it acts by translations.
Indeed a nontrivial linear elliptic automorphism has a
fixed point and its quotient is rational by
Riemann--Hurwitz. A rational general fibre on the
quotient would make every positive pluricanonical
section vanish by adjunction, contrary to $\kappa(M)=1$.
The finite translations form a fixed subgroup of the
constant elliptic curve, since points of bounded order
are discrete. We must retain the action at multiple
fibres when taking the quotient by $H$.

\par\medskip\noindent\textbf{The quotient by fibre translations.}
At a full fibre $mF$ of $q$, the actual normalized root
chart is
\[
 (F'\times\Delta_z)/\langle\delta\rangle,
 \qquad \delta(v,z)=(v+\tau,\zeta_m z),
 \qquad \operatorname{ord}(\tau)=m,
 \qquad t=z^m.                                        \tag{C.EQ.3}
\]
The root family is a smooth fixed-elliptic family.
The cyclic deck action is a translation on every fibre,
and its free action on the central fibre gives exactly
the indicated order of $\tau$. This is the normalized
root comparison proved in UN22 and UN29, with the
constant-translation gauge written in (C.KX.4).

Lift a base-fixing element of $H$ to this chart.
Its multiplier on $z$ is a power of $\zeta_m$;
composition with that inverse power of $\delta$
makes it fix $z$. The appropriate finite power
is a deck transformation fixing $z$, and hence
the identity. Its fibre action is therefore a
translation by a finite-order constant point.
The lifts so chosen generate a finite translation
subgroup $T\subset F'$, commuting with $\delta$.
First quotient by $T$. Write $\bar F=F'/T$ and
$m'=\operatorname{ord}(\bar\tau)$, so $m'\mid m$.
The remaining action is
$(\bar v,z)\mapsto(\bar v+\bar\tau,\zeta_m z)$.
Its $m'$-th power is a pure rotation of the disc.
Quotienting that subgroup replaces $z$ by
$w=z^{m/m'}$; the remaining free action is
\[
 (\bar v,w)\longmapsto(\bar v+\bar\tau,\zeta_{m'}w),
 \qquad t=w^{m'}.                                    \tag{C.EQ.4}
\]
This proves directly that $E_1=E/H$ is smooth near
the fibre and has the full smooth multiple fibre
$m'(\bar F/\langle\bar\tau\rangle)$. It is smooth
elsewhere by the same local translation quotient.
Thus $E_1\to B$ still has a fixed elliptic fibre,
identity integral monodromy, and only smooth multiples.
The remaining cyclic group $Q=G/H$ acts faithfully
on $B$.

\par\medskip\noindent\textbf{The local action at a base branch value.}
Let a stabilizer of order $s$ act at $b\in B$, and
let $\beta$ be its linear action on the fixed elliptic
curve. The linear action is constant along the regular
base, because it belongs to a discrete finite automorphism
group. Suppose first that $\beta\ne1$.
We prove that the fibre of $E_1$ at $b$ has multiplicity
one; this prevents a hidden change of the singular-fibre
type in the quotient.

Use (C.EQ.3) for a putative multiplicity $m$.
The generator lifts to the normalized root chart,
with normal multiplier $\rho$ satisfying
$\rho^m=\zeta_s$, and affine fibre derivative $\beta$.
It commutes with the deck action: its conjugation
preserves the multiplier $\zeta_m$ on the root
coordinate, and hence preserves the corresponding
deck generator. Restriction to the central elliptic
curve gives $\beta\tau=\tau$. The finite-order
endomorphism $\beta\ne1$ satisfies
$1+\beta+\cdots+\beta^{s-1}=0$. Applying it to
$\tau$ yields $s\tau=0$, so $m\mid s$.

The induced affine map on the reduced elliptic
fibre has a fixed point, because $\beta-1$ is a
surjective elliptic isogeny. Choose a lift of that
point, and compose the lifted generator with a
deck transformation to make it fix the chosen
lift. Its $s$-th power is a deck transformation
fixing a point of the central elliptic curve;
the deck action there is free, so that power is
the identity. Therefore $\rho^s=1$.
Together with $\rho^m$ having order $s$, this
gives $\gcd(m,s)=1$. Since $m\mid s$, we have
\[
                 m=1\qquad(\beta\ne1).               \tag{C.EQ.5}
\]
This is the same explicit local comparison used
in UN32, and the argument did not use a canonical
form or a K3 hypothesis.

When $\beta=1$, every element with a fixed point
on a reduced elliptic fibre fixes that whole fibre:
an elliptic translation with a fixed point is the
identity. Its stabilizer acts only on the normal
coordinate. A finite holomorphic action with
identity derivative is locally the identity,
by the averaging linearization, so every nontrivial
stabilizer is a normal reflection. The quotient
surface is smooth and the reduced fibre remains
smooth. Equivalently, the translation and root
quotients (C.EQ.3)--(C.EQ.4) display that fibre
as a smooth multiple. Hence these values have
zero topological Euler contribution.

For $\beta\ne1$, (C.EQ.5) gives a smooth product
near the source fibre. An equivariant local section
through a fixed point removes the translation part
of the finite action. To construct it, lift its
translation cocycle to small holomorphic complex-valued
functions near the fixed point and average over the
finite group. The cocycle sum is zero because the
action has finite order and the lift vanishes at
the fixed point. In additive notation, for an operator
$P$ of order $s$ and a cocycle $c$ with
$\sum_{k=0}^{s-1}P^kc=0$, the function
$s^{-1}\sum_{k=0}^{s-1}kP^kc$ solves
$(P-1)b=c$. This supplies the claimed holomorphic
section and gives a product action.
Quotient the subgroup acting trivially on the
elliptic factor. It consists of pure normal
reflections. The remaining action has order
$d\in\{2,3,4,6\}$ and the exact form
\[
             (v,z)\longmapsto(\zeta_d^r v,\zeta_d z),
             \qquad r\in\{1,d-1\}.                   \tag{C.EQ.6}
\]
Translations of an origin have already been removed.

Here is an explicit comparison with the Kodaira
types, including the direction of the character.
Write the fixed elliptic curve as
$y^2=x^3+Ax+B$, so multiplication by $\zeta_d^r$
acts on $x,y$ by $\zeta_d^{-2r},\zeta_d^{-3r}$.
The invariant meromorphic functions
\[
          T=z^d,\qquad X=z^{2r}x,\qquad Y=z^{3r}y
                                                               \tag{C.EQ.7}
\]
generate the quotient function field. They give
the following minimal Weierstrass equations:
\[
\begin{array}{c|c|c|c}
d&r&\text{equation}&\text{Kodaira type}\\ \hline
2&1&Y^2=X^3+AT^2X+BT^3&\mathrm I_0^*\\
3&1&Y^2=X^3+BT^2&\mathrm{IV}\\
3&2&Y^2=X^3+BT^4&\mathrm{IV}^*\\
4&1&Y^2=X^3+ATX&\mathrm{III}\\
4&3&Y^2=X^3+AT^3X&\mathrm{III}^*\\
6&1&Y^2=X^3+BT&\mathrm{II}\\
6&5&Y^2=X^3+BT^5&\mathrm{II}^* .
\end{array}                                                   \tag{C.EQ.8}
\]
For $d=3,6$ the curve has $A=0$ and $B\ne0$;
for $d=4$ it has $B=0$ and $A\ne0$.
The equations are minimal, since their indicated
orders are below the Weierstrass rescaling threshold.
Their discriminant orders and the resolved fibre
types are precisely the corresponding entries of
\cite[Theorem~6.2 and Table~I]{Kodaira1963II}.
The map (C.EQ.7) is birational over the base
after taking the finite quotient. Resolving its
graph and contracting vertical exceptional curves
gives the same relatively minimal elliptic model;
therefore these are the fibre types of the actual
quotient. This verifies the reduction to that table,
including the two different types of each primitive
order. Their Euler numbers are $12r/d$.

\par\medskip\noindent\textbf{All faithful cyclic base actions.}
On $B=\PP^1$, a nontrivial cyclic action is conjugate
to $z\mapsto\zeta_nz$. Its only branch values are
zero and infinity, where positive local generators
are inverse. A nontrivial fibre character of order
$d$ consequently gives the two opposite exponents
in (C.EQ.8). The resulting pairs are
\[
 (\mathrm{II},\mathrm{II}^*),\quad
 (\mathrm{III},\mathrm{III}^*),\quad
 (\mathrm{IV},\mathrm{IV}^*),\quad
 (\mathrm I_0^*,\mathrm I_0^*).                         \tag{C.EQ.9}
\]
Every pair has Euler sum twelve. A trivial fibre
character gives only smooth multiples.

On an elliptic $B$, a finite cyclic group either
acts by translations, in which case it is free,
or has a fixed point. In the latter case translate
that point to the origin. The action is linear of
order $n=2,3,4,6$, as follows from its action on
the rank-two integral period lattice. The branch
indices of $B\to B/Q\simeq\PP^1$ are
\[
\begin{array}{c|c}
n&\text{branch indices}\\ \hline
2&(2,2,2,2)\\
3&(3,3,3)\\
4&(4,4,2)\\
6&(6,3,2).
\end{array}                                                   \tag{C.EQ.10}
\]
For completeness, the number fixed by a nontrivial
power is $|1-\zeta_n^k|^2$, the determinant of
that power minus one on the lattice. Removing
the points with larger stabilizer and dividing
by the corresponding orbit sizes gives exactly
(C.EQ.10). Each row also satisfies the degree-$n$
Riemann--Hurwitz equation.

Write the fibre character as $\zeta_n^a$,
$0\le a<n$. At a branch index $e$, its exponent
is $a$ modulo $e$. Exponent zero gives a smooth
multiple. Otherwise put $c=\gcd(a,e)$; quotienting
the normal reflection kernel gives the entry of
(C.EQ.8) with $d=e/c$ and $r=(a\bmod e)/c$.
Enumerating these actual characters gives
\[
\begin{array}{c|c|l}
n&a&\text{nonsmooth fibre configuration}\\ \hline
2&1&4\mathrm I_0^*\\
3&1&3\mathrm{IV}\\
3&2&3\mathrm{IV}^*\\
4&1&2\mathrm{III}+\mathrm I_0^*\\
4&2&2\mathrm I_0^*\\
4&3&2\mathrm{III}^*+\mathrm I_0^*\\
6&1&\mathrm{II}+\mathrm{IV}+\mathrm I_0^*\\
6&2,4&\mathrm{IV}+\mathrm{IV}^*\\
6&3&2\mathrm I_0^*\\
6&5&\mathrm{II}^*+\mathrm{IV}^*+\mathrm I_0^* .
\end{array}                                                   \tag{C.EQ.11}
\]
The zero character gives no nonsmooth fibre.
These finite counts, the local reflection reductions,
and the Euler sums are independently executed in
\path{code/root_cyclic_elliptic_types.py} and its
saved output.

We finish by identifying the coherent invariants
without adding any unproved occurrence assertion.
The singular values just listed are all the possible
ones: a free base orbit is locally unchanged, and
the source had only smooth multiple fibres.
A free cyclic action on an elliptic base also leaves
only smooth multiples. Thus positive Euler occurs
only over $B_0=\PP^1$ and with nontrivial fibre
linear character. The exact smooth-multiple comparison
of UN29 applies to $E_1$ and gives $p_g(E_1)=g(B)$.
The generically finite quotient comparison, resolved
by its actual graph, injects holomorphic two-forms
on $M$ into the invariant forms on $E_1$, as in
(C.PG.4)--(C.PG.5).
For $B=\PP^1$ this immediately gives $p_g(M)=0$.
For elliptic $B$, the unique two-form on $E_1$ is
the wedge of a base differential and the fibre
differential with fixed normalized periods. At a
smooth multiple fibre the same form extends with
the zero of order $m-1$ supplied by the root
coordinate; hence it is a global holomorphic form.
Its character is $\zeta_n^{a+1}$. Consequently
$p_g(M)=0$ whenever $a+1\not\equiv0\pmod n$.
Every Euler-twelve row of (C.EQ.11) has precisely
that nontrivial character. This argument uses
pullback through the actual graph and Hartogs
extension, and does not infer extension of a
form from numerical invariance alone.

On a relatively minimal elliptic surface the canonical
divisor is vertical. An actual meromorphic differential
on the base and a meromorphic relative differential
give such a divisor; restriction to a general compact
elliptic fibre has no zero or pole. The components
of the displayed reducible singular fibres are $(-2)$
rational curves. A type $\mathrm{II}$ fibre has arithmetic
genus one and square zero, and a smooth reduced multiple
fibre is an elliptic curve of square zero. Adjunction gives
$K_M\cdot C=0$ for each of them, and hence $K_M^2=0$.
Analytic Riemann--Roch gives
$12\chi(\OO_M)=e(M)$. Euler localization for the
elliptic map $p$ and (C.EQ.8)--(C.EQ.11) give
exactly the rows of (C.EQ.1), or zero. For Euler
twelve we proved $p_g=0$, so $\chi=1$ gives $q=0$.
For Euler twenty-four, UN28 gives $p_g\le1$, while
$\chi=2=1-q+p_g$ and $q\ge0$ force $p_g=1,q=0$.
In the zero case
the local list has no nonsmooth reduced fibre,
since every such entry has positive Euler number.
This proves all assertions while preserving every
smooth multiple fibre and every actual quotient map.
\end{proof}


% C.2.UN.25 from part_C2_enriques_filling.tex

